\section{From the hypersurface inequality to model comparison}
\label{sec:comparison}

Fix $b\in\{0,1\}$, assume $\Sec_M\le-b^2$, and let $n=m+1\ge4$.
We use the confined profile $I$ in a geodesic ball $B=B_{R_B}(o)$
from Section~\ref{sec:profile}.  A minimizing boundary has the
decomposition $\Gamma=\Sigma\cup S$ of
Proposition~\ref{prop:confined-regularity}: $\Gamma$ is compact,
$S\Subset B$ is closed, and $\Sigma=\Gamma\setminus S$ is a
boundaryless $C^{1,1}_{\mathrm{loc}}$ hypersurface.  Its area measure
records the entire perimeter, including obstacle contact:
\begin{equation}
\label{eq:regular-perimeter-measure}
 P(E;D)=\sigma(\Sigma\cap D)
 \qquad\text{for every Borel set }D\subset M.
\end{equation}
At a differentiability volume $v$, Lemma~\ref{lem:profile-curvature} gives
$0\le h\le I'(v)/m$ almost everywhere on the regular boundary. The constant Sobolev test
would therefore bound the profile derivative from below in terms of
$I(v)$. Although $\Gamma$ is compact, its regular part
$\Sigma=\Gamma\setminus S$ need not be compact, so compactness of
$\Gamma$ alone does not justify that test. We will approximate it by compactly
supported cutoffs whose Dirichlet energy tends to zero. This requires
area growth near $S$, which we obtain directly from constrained minimality.

\subsection{Area growth by deletion and volume repair}

Deleting a small ball removes boundary near a singular point and creates
at most the area of the distance sphere. A variation on a fixed free
patch will restore the lost volume. We first record the deletion estimate
for finite-perimeter sets, so neither operation presupposes regularity
at the singular point.

\begin{lemma}[Deleting a geodesic ball]\label{lem:bv-deletion}
Let $E$ have finite volume and perimeter in a complete smooth
Riemannian manifold, let $x\in M$, and let
$0<r<\operatorname{inj}(x)$. Then
\begin{equation}\label{eq:bv-deletion}
 P(E\setminus B_r(x))\le
 P(E;M\setminus\overline B_r(x))
       +\mathcal H_g^{n-1}(\partial B_r(x)).
\end{equation}
For all but countably many radii in any bounded interval,
$P(E;\partial B_r(x))=0$, so the perimeters on the two open sides
sum to $P(E)$.
\end{lemma}

\begin{proof}
For a bounded Lipschitz scalar $\varphi$, distributional testing gives
\begin{equation}\label{eq:bv-multiplier}
 D(\varphi\chi_E)=\varphi D\chi_E+
                  \chi_E\nabla\varphi\,d\vol_g.
\end{equation}
First prove this for smooth multipliers. On a compact neighborhood of
a test field's support, approximate $\varphi$ smoothly, uniformly and
in $W^{1,1}$. Uniform convergence controls the measure term and
$W^{1,1}$ convergence controls the second term, since $|\chi_E|\le1$.
This proves the displayed distributional identity.

Take
\[
 \varphi_\varepsilon(y)=
 \min\{1,\max\{0,(d(x,y)-r)/\varepsilon\}\},
 \qquad r+\varepsilon<\operatorname{inj}(x).
\]
Its derivative is supported in a compact annulus. The products
$\varphi_\varepsilon\chi_E$ converge in $L^1$ to
$\chi_{E\setminus B_r(x)}$, because the smooth distance sphere has
ambient volume zero. Lower semicontinuity and
\eqref{eq:bv-multiplier} give
\[
 \begin{split}
 P(E\setminus B_r(x))
 &\le\liminf_{\varepsilon\downarrow0}
 \left(\int\varphi_\varepsilon\,d|D\chi_E|
       +\frac{|E\cap(B_{r+\varepsilon}(x)\setminus B_r(x))|}
                    {\varepsilon}\right)\\
 &\le P(E;M\setminus\overline B_r(x))
       +\left.\frac d{ds}\right|_{s=r}|B_s(x)|.
 \end{split}
\]
Polar coordinates identify the last derivative with the area of the
distance sphere. Finally, disjoint spheres can carry positive mass
for a finite measure at only countably many radii.
\end{proof}

\begin{lemma}[Area growth at singular points]
\label{lem:area-growth}
Let $0<v<\vol(B)$, and let $E$ be a confined minimizing region of
volume $v$.  There are constants $C<\infty$ and $r_0>0$, depending on
$E$ and $B$, such that
\begin{equation}
\label{eq:singular-area-growth}
 \sigma\bigl(\Sigma\cap B_r(x)\bigr)\le Cr^m
 \qquad(x\in S,\ 0<r<r_0).
\end{equation}
\end{lemma}

\begin{proof}
Suppose $S\ne\varnothing$, since otherwise the assertion is empty.
By \eqref{eq:positive-free-perimeter} and
\eqref{eq:perimeter-area}, there is a free smooth boundary point.
Choose a neighborhood $W$ of that point with
$\overline W\subset B\setminus S$ so small that $\Gamma\cap W$ is
smooth.  Extend a nonnegative, nonzero normal variation of a smaller
patch to a smooth field $Z$ compactly supported in $W$.  With $N$ the
outward normal, this gives
\[
 a:=\int_{\Sigma\cap W}\langle Z,N\rangle\,d\sigma>0.
\]

Let $\phi_t$ be its flow and put
$G(t)=\vol(\phi_t(E))-\vol(E)$.  The volume Jacobian formula gives
$G'(0)=a$.  Choose $t_0>0$ so that $G'(t)\ge a/2$ for
$0\le t\le t_0$.  The flow preserves $W$ and $B$ and fixes
$M\setminus W$.  Whenever a finite-perimeter set $F$ agrees with $E$
almost everywhere on $W$, the volume Jacobian minus one is supported
in $W$, so
\begin{equation}
\label{eq:fixed-patch-volume}
 \vol(\phi_t(F))-\vol(F)=G(t).
\end{equation}
The tangential Jacobian formula and locality of the reduced boundary
likewise give
\[
 P(\phi_t(F))-P(F)
 =\int_{\partial^*E\cap W}
 \left(J_m\bigl(D\phi_t|_{T_x\partial^*E}\bigr)-1\right)
 \,d\mathcal H_g^m(x).
\]
On the compact support of the flow, the tangential Jacobians satisfy
$|J_m-1|\le Lt$ uniformly over all tangent $m$-planes for
$0\le t\le t_0$.  Hence
\begin{equation}
\label{eq:fixed-patch-perimeter}
 \bigl|P(\phi_t(F))-P(F)\bigr|\le L P(E;W)t.
\end{equation}
Every volume deficit $0\le\delta\le at_0/2$ can therefore be
repaired by a time $0\le t\le2\delta/a$, at perimeter cost at most
$C_R\delta$, where $C_R=2LP(E;W)/a$.  The bound is independent of
$F$ as long as $F=E$ almost everywhere on $W$.

Compactness of $S$ and its disjointness from
$\overline W\cup\partial B$ give $\rho_*>0$ such that
\[
 B_{2\rho_*}(x)\subset B\setminus\overline W
 \qquad(x\in S).
\]
After decreasing $\rho_*$, smoothness of the ambient metric on a
compact neighborhood of $S$ gives uniform constants $C_V,C_A$ with
\begin{equation}
\label{eq:small-ambient-balls}
 \vol(B_r(x))\le C_Vr^n,
 \qquad
 \mathcal H_g^m(\partial B_r(x))\le C_Ar^m
 \quad(x\in S,\ 0<r<\rho_*).
\end{equation}
Decrease $\rho_*$ again so that $C_V\rho_*^n\le at_0/2$.

Fix $x\in S$.  For almost every $r\in(0,\rho_*)$, one has
$P(E;\partial B_r(x))=0$. At such radii set $F=E\setminus B_r(x)$;
Lemma~\ref{lem:bv-deletion}, applied to $E$, gives
\begin{equation}
\label{eq:deleted-perimeter}
 P(F)\le P(E;M\setminus\overline{B_r(x)})
       +\mathcal H_g^m(\partial B_r(x)).
\end{equation}
The volume deficit is
$\delta=\vol(E\cap B_r(x))\le C_Vr^n$, and $F=E$ on $W$.
Use \eqref{eq:fixed-patch-volume} to repair this deficit by a flow
time $t\le2\delta/a$.  The repaired set remains in $B$ and has
volume $v$.  Minimality, \eqref{eq:fixed-patch-perimeter}, and
\eqref{eq:deleted-perimeter} imply
\[
 \begin{split}
 P(E;B_r(x))
 &\le\mathcal H_g^m(\partial B_r(x))+C_R\delta\\
 &\le C_Ar^m+C_RC_Vr^n\le C_0r^m,
 \end{split}
\]
where $C_0=C_A+C_RC_V\rho_*$ because $n=m+1$.
The constants are uniform in $x\in S$.

For an arbitrary $0<r<\rho_*/2$, choose a good radius
$\rho\in(r,2r)$.  Monotonicity of the perimeter measure gives
\[
 P(E;B_r(x))\le P(E;B_\rho(x))\le2^m C_0r^m.
\]
Together with \eqref{eq:regular-perimeter-measure}, this proves the
claim with $r_0=\rho_*/2$.
\end{proof}

\subsection{Removing the singular set from the constant test}

The area bound turns a cutoff gradient of order $r^{-1}$ on a ball of
radius $r$ into an energy cost of order $r^{m-2}$. The small Hausdorff
dimension of the singular set will make the sum of these costs tend to
zero. We isolate the cutoff construction because the same conclusion
will later apply to the boundary of an equality set.

\begin{lemma}[Cutoffs around a set of zero capacity]\label{lem:capacity-cutoffs}
Let $\Gamma$ be a compact subset of a smooth Riemannian manifold, and
let $S\subset\Gamma$ be closed. Suppose that
$\Sigma=\Gamma\setminus S$ is an embedded, boundaryless, locally
$C^{1,1}$ hypersurface of dimension $m\ge3$ and finite area.
Assume $\mathcal H^{m-2}(S)=0$ and that, for some $C,r_0>0$,
\[
 \sigma(\Sigma\cap B_r(x))\le Cr^m
 \qquad(x\in S,\ 0<r<r_0).
\]
There are $\eta_j\in\Lip_c(\Sigma)$ such that
\begin{equation}\label{eq:singular-cutoff-properties}
 0\le\eta_j\le1,\qquad \eta_j(x)\longrightarrow1\quad(x\in\Sigma),
 \qquad \int_\Sigma|\nabla_\Sigma\eta_j|^2\,d\sigma\longrightarrow0.
\end{equation}
If $S=\varnothing$, the choice $\eta_j=1$ works in every dimension.
\end{lemma}

\begin{proof}
If $S=\varnothing$, then $\Sigma=\Gamma$ is compact and we take
$\eta_j=1$, including when $\Sigma$ has several components.
Otherwise, use $\mathcal H^{m-2}(S)=0$ as follows.
For each
positive integer $j$, choose a finite covering
\begin{equation}
\label{eq:singular-cover}
 \begin{gathered}
 S\subset\bigcup_{i=1}^{N_j}B_{r_{ji}}(x_{ji}),
 \qquad x_{ji}\in S,\\
 0<r_{ji}<\min(2^{-j},r_0/4),
 \qquad\sum_{i=1}^{N_j}r_{ji}^{m-2}<2^{-j}.
 \end{gathered}
\end{equation}
Indeed, Hausdorff nullity first gives a countable open-ball cover
with arbitrarily small radii and sum.  Recenter the balls at points
of $S$ with a fixed enlargement, starting with tolerances small enough
to absorb this enlargement, and take a finite subcover by compactness.

Put $\theta(t)=\min\{1,\max\{0,t-1\}\}$ for $t\ge0$, and define
\[
 \eta_{ji}(x)=\theta\!\left(\frac{d(x,x_{ji})}{r_{ji}}\right),
 \qquad
 \eta_j=\min_{1\le i\le N_j}\eta_{ji}\big|_\Sigma.
\]
Each ambient cutoff vanishes on $B_{r_{ji}}(x_{ji})$, is one
outside $B_{2r_{ji}}(x_{ji})$, and has Lipschitz constant at most
$r_{ji}^{-1}$.  A finite minimum is Lipschitz, and almost everywhere
\[
 |\nabla_\Sigma\eta_j|^2
 \le\sum_{i=1}^{N_j}|\nabla_\Sigma\eta_{ji}|^2.
\]
There is no restriction on the overlap of the balls.  Applying the assumed area bound at their doubled radii gives
\begin{equation}
\label{eq:singular-cutoff-energy}
 \begin{split}
 \int_\Sigma|\nabla_\Sigma\eta_j|^2\,d\sigma
 &\le\sum_i r_{ji}^{-2}
       \sigma(\Sigma\cap B_{2r_{ji}}(x_{ji}))\\
 &\le2^m C\sum_i r_{ji}^{m-2}<2^m C\,2^{-j}.
 \end{split}
\end{equation}

The open balls in \eqref{eq:singular-cover} form a neighborhood of
$S$ on which $\eta_j$ vanishes.  Its support is therefore contained
in a compact subset of $\Gamma\setminus S$.  The embedded graph
charts identify the subspace and manifold topologies on $\Sigma$,
so this support is compact in $\Sigma$.  For each fixed
$x\in\Sigma$, closedness of $S$ gives $d(x,S)>0$.  As the maximum
radius tends to zero, the doubled balls eventually miss $x$, and
$\eta_j(x)=1$.  This proves
\eqref{eq:singular-cutoff-properties}, without assuming completeness
of $\Sigma$ or finiteness of its components.

\end{proof}

\begin{lemma}[The constant test]
\label{lem:constant-test}
Let $v\in(0,\vol(B))$ be a differentiability point of $I$, and let
$E$ minimize the confined perimeter at volume $v$.  Then
\begin{equation}
\label{eq:constant-sobolev}
 m\int_\Sigma(h^2-b^2)\,d\sigma
 \ge Y_m I(v)^{(m-2)/m}.
\end{equation}
\end{lemma}

\begin{proof}
Lemma~\ref{lem:profile-curvature} gives
$0\le h\le h_0:=I'(v)/m<\infty$ almost everywhere on $\Sigma$.
If $S\ne\varnothing$, the singular-dimension bound gives
\[
 \dim_{\mathcal H}S\le n-8=m-7<m-2,
 \qquad \mathcal H^{m-2}(S)=0.
\]
Lemma~\ref{lem:area-growth} and the compact embedded boundary structure
therefore verify the hypotheses of Lemma~\ref{lem:capacity-cutoffs}.
That lemma supplies cutoffs with \eqref{eq:singular-cutoff-properties};
when $S=\varnothing$ take $\eta_j=1$.

Apply Theorem~\ref{thm:extrinsic-sobolev} to these cutoffs:
\[
 c\int_\Sigma|\nabla_\Sigma\eta_j|^2\,d\sigma
 +m\int_\Sigma(h^2-b^2)\eta_j^2\,d\sigma
 \ge Y_m\left(\int_\Sigma\eta_j^q\,d\sigma\right)^{2/q}.
\]
The area is $I(v)<\infty$, and
$|h^2-b^2|\le h_0^2+b^2$ almost everywhere.  Dominated convergence
and \eqref{eq:singular-cutoff-properties} prove
\eqref{eq:constant-sobolev}, since $2/q=(m-2)/m$.
The same construction applies in every dimension under consideration:
when $4\le n\le7$ the singular set is empty, and when $n=8$ its
possible infinite zero-dimensional compact set still admits the
covers in \eqref{eq:singular-cover}.
\end{proof}

\subsection{The profile differential inequality and its integral}

The singular set has now been removed from the constant test without
assuming completeness of the regular locus. The curvature bound from
the confined profile can therefore be inserted in the sharp Sobolev
inequality. This is the step that reduces the geometric problem to one
scalar differential inequality.

\begin{proposition}[Profile differential inequality]
\label{prop:profile-ode}
For almost every $v\in(0,\vol(B))$,
\begin{equation}
\label{eq:profile-differential}
 I'(v)\ge m\sqrt{b^2+\left(\frac{s_m}{I(v)}\right)^{2/m}}.
\end{equation}
\end{proposition}

\begin{proof}
The profile is locally Lipschitz by
Lemma~\ref{lem:profile-lipschitz}, hence differentiable almost
everywhere.  At a differentiability volume choose a minimizer.
Lemma~\ref{lem:profile-curvature} gives
$h_0=I'(v)/m>0$ and $h^2\le h_0^2$ almost everywhere on its regular
boundary.  Lemma~\ref{lem:constant-test} therefore gives
\[
 m(h_0^2-b^2)I(v)\ge Y_m I(v)^{(m-2)/m}.
\]
Since $I(v)>0$ and $Y_m=m s_m^{2/m}$, division yields
\[
 h_0^2\ge b^2+\left(\frac{s_m}{I(v)}\right)^{2/m}.
\]
Taking the positive square root proves
\eqref{eq:profile-differential}.
\end{proof}

\begin{theorem}[Comparison in the normalized curvatures]
\label{thm:normalized-comparison}
Let $n\ge4$ and $b\in\{0,1\}$.  If $M^n$ is complete, simply
connected and smooth, with $\Sec_M\le-b^2$, then every bounded measurable
set $E\subset M$ of finite perimeter and positive volume satisfies
\[
 P(E)\ge\mathcal I_b(\vol(E)).
\]
\end{theorem}

\begin{proof}
Choose radii $0<R_0<R$ such that $E\subset B_{R_0}(o)$ up to a
null set, and put $B=B_R(o)$. The annulus
$B_R(o)\setminus\overline{B_{R_0}(o)}$ has positive volume, so
$\vol(B)>\vol(E)$.  Recall the model area and volume functions
\[
 \mathcal A_b(r)=s_m\sn_b(r)^m,
 \qquad
 \mathcal V_b(r)=s_m\int_0^r\sn_b(t)^m\,dt.
\]
Both are strictly increasing from zero to infinity.  For $A>0$ put
\begin{equation}
\label{eq:inverse-model-profile}
 R_b(A)=\mathcal A_b^{-1}(A),
 \qquad
 \mathcal W_b(A)=\mathcal V_b(R_b(A)).
\end{equation}
Thus $\mathcal W_b$ is strictly increasing.  The identities
$\sn_b'=\cs_b$ and $\cs_b^2-b^2\sn_b^2=1$ give
\begin{equation}
\label{eq:inverse-profile-derivative}
 \begin{split}
 \mathcal W_b'(A)
 &=\frac{\sn_b(R_b(A))}{m\cs_b(R_b(A))}\\
 &=\frac{1}{m\sqrt{b^2+(s_m/A)^{2/m}}}.
 \end{split}
\end{equation}

On every interval $[\varepsilon,v]\subset(0,\vol(B))$, the
positive Lipschitz function $I$ has image in a compact subset of
$(0,\infty)$, where $\mathcal W_b$ is continuously differentiable.
Thus $\mathcal W_b\circ I$ is absolutely continuous there.
Proposition~\ref{prop:profile-ode}, the chain rule, and
\eqref{eq:inverse-profile-derivative} imply
\[
 (\mathcal W_b\circ I)'(s)
 =\mathcal W_b'(I(s))I'(s)\ge1
 \quad\text{for almost every }s\in[\varepsilon,v].
\]
Integrating yields
\[
 \mathcal W_b(I(v))
 \ge\mathcal W_b(I(\varepsilon))+v-\varepsilon
 \ge v-\varepsilon.
\]
Letting $\varepsilon\downarrow0$, we obtain
\begin{equation}
\label{eq:integrated-profile}
 \mathcal W_b(I(v))\ge v
 \qquad(0<v<\vol(B)).
\end{equation}
This endpoint passage uses only nonnegativity of
$\mathcal W_b(I(\varepsilon))$.

Set $V=\vol(E)$. The choice of $B$ gives $0<V<\vol(B)$, and $E$
is an admissible competitor in \eqref{eq:confined-profile}. Consequently
\[
 P(E)\ge I(V).
\]
This step uses the full ambient perimeter of $E$, with no regularity
assumption on its boundary.
If $r>0$ is determined by $\mathcal V_b(r)=V$, then
$\mathcal W_b(\mathcal A_b(r))=V$.
Equation~\eqref{eq:integrated-profile} and strict monotonicity of
$\mathcal W_b$ give $I(V)\ge\mathcal A_b(r)=\mathcal I_b(V)$,
which proves the assertion.
\end{proof}
